% !Mode::"TeX:UTF-8"

% -------------------- Information --------------------

\newcommand{\TITLE}{幻灯片模板}
\newcommand{\AUTHOR}{Jason}
\newcommand{\SUBJECT}{幻灯片}
\newcommand{\KEYWORDS}{}
\newcommand{\INSTITUTE}{同济大学数学科学学院}

% -------------------- Packages --------------------

\documentclass[xcolor=dvipsnames]{ctexbeamer}
\usepackage{amsmath}
%\usepackage{amsthm} % loaded by beamer.
\usepackage{amssymb}
\usepackage{commath} % abs, norm
\usepackage[mathscr]{euscript}
\usepackage{float} % 你们这帮float给我乖乖听话 HHHHHHHHHHH.
\usepackage{graphicx}
%\usepackage{hyperref} % loaded by beamer.
\usepackage{mathtools} % \xleftrightarrow.
\usepackage{tikz}
%\usepackage{xcolor} % loaded by beamer.
\usepackage[printwatermark]{xwatermark} % Foreground Watermarks.
\usepackage[all, cmtip]{xy}

% -------------------- Settings --------------------

% Title

\title{\TITLE}
\author{\AUTHOR}
\date{\today}
\institute{\INSTITUTE}
\titlegraphic{\includegraphics[width=0.1\paperwidth]{wallpaper.jpg}}

% Theorem Environments

\let\note\undefined
\newtheorem{note}{注}[section]

% Package: beamer

\transpush<2>[duration=2]
\usetheme{Berlin}

% Package: ctex

\setCJKfamilyfont{fzstk}{FZShuTi} % 方正舒体
\newcommand{\fzstk}{\CJKfamily{fzstk}}

% Package: graphicx

\graphicspath{{resources/}} % 图像文件目录

% Package: hyperref

\hypersetup{
    linktoc             =   all,
    colorlinks          =   true,
    linkcolor           =   cyan,
    anchorcolor         =   black,
    citecolor           =   green,
    filecolor           =   cyan,
    menucolor           =   red,
    runcolor            =   filecolor,
    urlcolor            =   magenta,
    pdftitle           	=   {\TITLE},
    pdfauthor          	=   {\AUTHOR},
    pdfsubject         	=   {\SUBJECT},
    pdfcreator			=	{Visual Studio Code},
    pdfproducer			=	{XeLaTeX with documentclass ctexbeamer},
    pdfkeywords        	=   {\KEYWORDS},
    bookmarksnumbered   =   true,
    pdfstartview        =   Fit,
    pdfpagelayout       =   OneColumn,
    pdfpagemode			=   FullScreen
}

}

% Package: xwatermark

\newsavebox\mybox
\savebox\mybox{\tikz[color=cyan, opacity=0.2]\node{\fzstk\SUBJECT};}
\newwatermark*[
    allpages,
    angle=45,
    scale=6,
    xpos=-20,
    ypos=15
]{\usebox\mybox}

% -------------------- General new commands --------------------

\DeclareMathAlphabet{\mathsfsl}{OT1}{cmss}{m}{sl}

\newcommand{\diff}{\mathop{}\!\mathrm{d}}
\newcommand{\matr}[1]{\ensuremath{\mathsfsl{#1}}} % italic sans serif
\newcommand{\me}{\mathrm{e}}
\newcommand{\mi}{\mathrm{i}}
\newcommand{\restrict[1]}{\raisebox{-.5ex}{$|$}_{#1}}

% -------------------- Specific new commands --------------------

\newcommand{\mahjong}{\mathscr{M}}
\newcommand{\point}{\mathrm{point}}
\newcommand{\Hand}{\mathscr{H}}
\newcommand{\hand}{h}
\newcommand{\base}{\omega}
\newcommand{\prob}{\mathrm{Pr}}

% -------------------- Document --------------------

\begin{document}

    % -------------------- Title Page --------------------

    \begin{frame}
        \titlepage
    \end{frame}
    
    % -------------------- Contents --------------------
    
    \begin{frame}
        \frametitle{目录}
        \tableofcontents
    \end{frame}

    % -------------------- Body --------------------
    
    \section{背景}

    \begin{frame}
        \frametitle{背景}
        \begin{itemize}
            \item 麻将运动在我国有广泛的群众基础, 并且在2017年4月成为了继桥牌,
                国际象棋, 围棋, 中国象棋和国际跳棋以后第六个世界智力运动会项目.
            \item 不过从数学角度研究麻将的论文很少,
                在arXiv上搜索`Mahjong'关键词仅有两条相关结果.
            \item 其中一篇
                "Mathematical aspects of the combinatorial game Mahjong"
                [arXiv:1707.07345v2]研究麻将牌中的组合问题,
                与麻将博弈无关.
            \item 另一篇"Let's Play Mahjong!"[arXiv:1903.03294v1]
                花了主要篇幅定义并计算了牌型的好坏程度,
                而只在一章里面提到了一点点具体决策的方法.
        \end{itemize}
    \end{frame}

    \section{基本定义}

    \begin{frame}
        \frametitle{博弈的基本要素}
        博弈的五要素分别为局中人, 策略集, 收益, 信息和行动顺序.
        我们将从这五个要素来介绍本麻将博弈的基本要素.
    \end{frame}

    \begin{frame}
        \frametitle{局中人}
        \begin{definition}[局中人]
            局中人就是在麻将博弈中独自决策的个体, 最后获得相应结果.
        \end{definition}
        
        \begin{block}{假设}
            \begin{itemize}
                \item 假设局中人做的决策与局中人追求的目的是一致的.
                \item 假设局中人的目的是最大化收益的期望值.
                \item {[\alert{完全理性假设}]}假设局中人有一切博弈所需的分析能力,
                    并且在行为上不会犯失误等任何错误.
            \end{itemize}
        \end{block}
    \end{frame}

    \begin{frame}
        \frametitle{策略}
        \begin{definition}[策略]
            一个策略是一个局中人的完整的相机抉择的行动计划,
            在本博弈问题中即局中人对鸣牌(吃碰杠)和(胡)切牌的选择.
        \end{definition}
        \begin{definition}[策略集]
            一个局中人的所有的策略的全体称为这个局中人的策略集.
        \end{definition}
    \end{frame}

    \begin{frame}
        \frametitle{收益}
        \begin{definition}[收益]
            收益是策略组合的函数, 在本博弈中即为一局博弈结束时的点数得失.
        \end{definition}
    \end{frame}

    \begin{frame}
        \frametitle{信息}
        \begin{definition}[信息]
            信息是与博弈有关的消息和知识.
        \end{definition}
        \begin{definition}[消息]
            消息, 之后我们称为消息库$\base$, 是局中人之间的共同知识.
            在本博弈中为每个局中人切出的牌和切牌顺序, 以及鸣(吃碰杠)的牌.
        \end{definition}
        \begin{note}[消息库]
            消息库反映了其它局中人手牌的可能组合以及牌堆中牌的可能组合.
        \end{note}
    \end{frame}

    \begin{frame}
        \frametitle{信息(续)}
        \begin{definition}[知识]
            知识在本博弈中指的是博弈的相关定义以及麻将的规则等.
        \end{definition}
        \begin{note}[私人信息]
            所有局中人的私人信息有且仅有自己当前手牌.
        \end{note}
    \end{frame}
    
    \begin{frame}
        \frametitle{行动顺序}

    \end{frame}

    \begin{frame}
        \frametitle{麻将规则的定义}
        \begin{definition}[麻将规则]
            麻将规则由麻将牌全体和每种14张手牌的点数构成, 我们定义了如下记号:
            \begin{itemize}
                \item 麻将牌全体为广义集合$\mahjong$,
                    合法手牌全集为$\Hand\subset\mahjong$,
                    且有$\abs{\hand}=14, \forall{\hand\in\Hand}$.
                \item 手牌点数为映射$\point: \Hand\rightarrow\mathbb{N}$.
                \item 和牌手牌(winning hands)全体$\mathscr{W}=\{\point>0\}$.
            \end{itemize}
        \end{definition}
    \end{frame}

    \section{纯进攻}

    \begin{frame}
        \frametitle{纯进攻游戏引言}
        \begin{itemize}
            \item 麻将博弈中一局中人只有和牌才可能得到可观的正收益,
                因此尽快和牌的能力也就是做牌的能力十分重要.
            \item 所以我们将引入纯进攻游戏,
                通过将消息库中其余局中人的私人消息转化为牌山的消息,
                将博弈问题分解为最优化问题和计算牌山可能组合概率的问题.
        \end{itemize}
    \end{frame}

    \begin{frame}
        \frametitle{纯进攻游戏定义}
        \begin{definition}[纯进攻游戏]
            该游戏仅有一位玩家, 该玩家手中有十四张牌$T$,
            他的消息库$\base$为牌堆中牌的所有可能组合及其概率.
            每回合该玩家打出一张牌$i$并随机摸进一张牌$t$, 消息库也随之更新,
            更新函数为$\base\mapsto\mathrm{renew}_t{\base}$.
            给定限定回合数$k$,
            若该玩家在$k$回合内和牌, 则收益记为该和牌型的点数,
            若该玩家未能在$k$回合内和牌, 则收益为0.
        \end{definition}
        记消息库为$\base$时
        局中人选择切出手牌$i$在$k$回合内完成点数为$p$的和牌的概率为
        $\mathrm{val}_{k}^{p}(\hand, \base, i)$,
        则纯进攻游戏优化的目标即为期望
        \[\sum_{p}{p\cdot\mathrm{val}_{k}^{p}(\hand, \base, i)}.\]
    \end{frame}

    \begin{frame}
        \frametitle{求解纯进攻游戏思路}
        定义
        \[\mathrm{val}_{k}^{p}(\hand, \base)=
        \max_{i\in\mathbb{Z}_{14}}{\mathrm{val}_{k}^{p}(\hand, \base, i)}.\]
        那么对于每一个$p$, 都应有
        \[\mathrm{val}_{k}^{p}(\hand, \base, i)=
        \sum_{t}{\Pr_{\base}}(t)\mathrm{val}_{k-1}^{p}(\hand[i/t],
        \mathrm{renew}_{t}^{-1}{\base}),\]
        其中$\Pr_{\base}(t)$为摸进$t$的概率,
        $T[i/t]$为$T$打出$i$并摸进$t$后的手牌.
        于是我们可递归求解$\mathrm{val}_{k}^{p}(\hand, \base, i)$,
        从而解得纯进攻游戏各决策的期望收益.
    \end{frame}
\end{document}